\documentclass[11pt, a4paper]{article}
\usepackage[margin=2.5cm]{geometry}
\usepackage{fontenc}
\usepackage{inputenc}
\usepackage{microtype}
\usepackage{parskip}
\usepackage{titlesec}
\usepackage{enumitem}
\usepackage{booktabs}
\usepackage{array}
\usepackage{xcolor}
\usepackage{mdframed}
\usepackage{listings}
\usepackage{hyperref}

\hypersetup{
    colorlinks=true,
    linkcolor=blue!60!black,
    urlcolor=blue!60!black,
}

\definecolor{codebg}{gray}{0.95}
\definecolor{alertbg}{RGB}{255, 248, 235}
\definecolor{alertborder}{RGB}{230, 150, 0}
\definecolor{okbg}{RGB}{240, 250, 240}
\definecolor{okborder}{RGB}{60, 160, 60}
\definecolor{highred}{RGB}{180, 30, 30}
\definecolor{medorange}{RGB}{180, 100, 0}
\definecolor{lowgreen}{RGB}{40, 130, 40}

\lstset{
    backgroundcolor=\color{codebg},
    basicstyle=\ttfamily\small,
    breaklines=true,
    frame=single,
    framesep=4pt,
    rulecolor=\color{gray!50},
    xleftmargin=8pt,
    xrightmargin=8pt,
    aboveskip=6pt,
    belowskip=6pt,
}

\titleformat{\section}{\large\bfseries}{}{0em}{}[\titlerule]
\titleformat{\subsection}{\normalsize\bfseries}{}{0em}{}

\newmdenv[
    backgroundcolor=alertbg,
    linecolor=alertborder,
    linewidth=2pt,
    leftline=true,
    rightline=false,
    topline=false,
    bottomline=false,
    innerleftmargin=10pt,
    innerrightmargin=8pt,
    innertopmargin=6pt,
    innerbottommargin=6pt,
]{alertbox}

\newmdenv[
    backgroundcolor=okbg,
    linecolor=okborder,
    linewidth=2pt,
    leftline=true,
    rightline=false,
    topline=false,
    bottomline=false,
    innerleftmargin=10pt,
    innerrightmargin=8pt,
    innertopmargin=6pt,
    innerbottommargin=6pt,
]{okbox}

\newcommand{\priority}[2]{%
    \textbf{\textcolor{#1}{#2}}%
}

\begin{document}

\begin{center}
    {\LARGE\bfseries AVU Sync Code Review}\\[4pt]
    {\large What Needs Human Attention Before Merging}\\[8pt]
    {\small Branch: \texttt{feature/avu-sync} \quad|\quad Generated: 2026-04-01}
\end{center}

\vspace{4pt}
\hrule
\vspace{12pt}

\section{Context}

The \texttt{feature/avu-sync} branch (AI-written) added \texttt{kando/sync/} --- a module for
automatically syncing iRODS AVU metadata to CKAN on a scheduled basis. The metadata mapping in
\texttt{kando/sync/mapping.py} is a refactor of the production logic in
\texttt{kando/helpers/migration.py}. This document records what has been verified correct and
what requires careful human review before the branch is merged or deployed.

\section{Already Verified Correct (matches \texttt{de.py})}

\begin{okbox}
\texttt{kando/sync/irods\_client.py} was cross-checked line-by-line against the production
\texttt{kando/de.py}. All API patterns match.
\end{okbox}

\vspace{6pt}

\begin{itemize}[nosep, leftmargin=1.5em]
    \item \textbf{Auth endpoint:} \texttt{/token/keycloak} --- identical
    \item \textbf{Directory listing:} \texttt{/secured/filesystem/directory?path=...} returns \texttt{\{"folders": [...]\}} --- identical
    \item \textbf{Metadata endpoint:} \texttt{/filesystem/\{id\}/metadata} --- identical
    \item \textbf{AVU parsing:} \texttt{avu['attr']} / \texttt{avu['value']} loop --- identical
    \item \textbf{Date format:} \texttt{strftime('\%Y-\%m-\%d \%H:\%M:\%S')} --- identical
    \item \textbf{Env vars:} \texttt{DE\_USERNAME}, \texttt{DE\_PASSWORD}, \texttt{TERRAIN\_URL} --- already used by \texttt{de.py}
\end{itemize}

\vspace{6pt}

The metadata field mapping in \texttt{mapping.py} is functionally equivalent to \texttt{migration.py}:
same key fallback order, same license logic, same tag normalization, same \texttt{dont\_include} set.
The AVU sync is actually cleaner in two minor ways: \texttt{get\_description()} handles list values by
joining them (instead of returning a raw Python list), and \texttt{get\_extras()} is more defensive
about missing date fields.

\section{Issues Requiring Human Review}

\subsection{\priority{highred}{[HIGH]}\quad \texttt{mapping.py:157} --- \texttt{create\_citation()} called unconditionally}

\begin{alertbox}
\textbf{Risk:} Datasets without an author or publication year will crash silently during sync.
\end{alertbox}

\texttt{get\_extras()} always calls \texttt{create\_citation()}, which in turn calls \texttt{get\_author()}
and \texttt{get\_publication\_year()}. Both raise \texttt{KeyError} if their respective fields are absent
from the AVU metadata.

\texttt{sync\_avu.py} only guards against a missing title (lines 104--109):

\begin{lstlisting}[language=Python]
try:
    get_title(avus)
except KeyError:
    logger.warning("Skipping %s: no title in metadata", dirname)
    stats["skipped"] += 1
    continue
\end{lstlisting}

A dataset with a title but no \texttt{datacite.creator}/\texttt{creator} or no publication year will
crash inside \texttt{map\_avus\_to\_ckan()} at the \texttt{get\_extras()} call. The outer
\texttt{except Exception} on line~179 catches it and logs an error, but the message
(\texttt{"No creator found in metadata"}) gives no context about which dataset failed or why.

\textbf{What to verify:} Do all curated datasets on prod have author and publication year?
Run a check against the Terrain API before deploying. Consider either:
\begin{itemize}[nosep, leftmargin=2em]
    \item Adding the same \texttt{try/except KeyError} guard for author and year in \texttt{sync\_avu.py}, or
    \item Wrapping \texttt{create\_citation()} in a \texttt{try/except} inside \texttt{get\_extras()} and falling back to an empty string.
\end{itemize}

\subsection{\priority{medorange}{[MEDIUM]}\quad \texttt{state.py:12} --- Default state file path}

\begin{lstlisting}[language=Python]
DEFAULT_STATE_PATH = os.path.join(os.path.dirname(__file__), "..", "sync_state.json")
\end{lstlisting}

This resolves to \texttt{kando/sync\_state.json}. The file accumulates CKAN dataset IDs and
iRODS timestamps for every curated dataset and will grow over time.

\textbf{What to verify:} Is \texttt{kando/sync\_state.json} listed in \texttt{.gitignore}?
It should not be committed to the repository.

\subsection{\priority{medorange}{[MEDIUM]}\quad \texttt{sync\_avu.py:120--172} --- Create/update orchestration logic}

This is the most complex block in the AI-written code. Read it carefully end-to-end.

\textbf{Create path (lines 121--152):}
\begin{itemize}[nosep, leftmargin=2em]
    \item If a dataset name already exists in CKAN, it updates rather than creates. Correct.
    \item If a name collision occurs on \texttt{package\_create}, it retries with a DOI suffix appended to the name. Verify the suffix truncation logic (\texttt{name[:80] + "\_" + suffix[:19]}) won't produce invalid CKAN names.
    \item \texttt{stats["created"] += 1} on line~152 is reached even when the ``already exists, updating'' path (lines 123--127) was taken. The created count will be inflated for re-syncs.
\end{itemize}

\textbf{Update path (lines 153--172):}
\begin{itemize}[nosep, leftmargin=2em]
    \item If \texttt{state.get\_ckan\_id()} returns \texttt{None} (lost state), the code tries to find the dataset by name, then falls back to creating a new one.
    \item \textbf{Walk this path manually} to confirm it cannot create a duplicate dataset when the dataset already exists in CKAN under the same name.
    \item \texttt{stats["updated"] += 1} on line~172 is always hit, even if a new dataset was created. Minor, but misleading.
\end{itemize}

\textbf{State tracking:} \texttt{state.mark\_synced()} is called at line~177 with a \texttt{ckan\_id}
guard. This covers all successful paths correctly.

\subsection{\priority{lowgreen}{[LOW]}\quad \texttt{example.env} --- Missing env vars for the sync}

\texttt{DE\_USERNAME} and \texttt{DE\_PASSWORD} are required by both the existing \texttt{de.py} and
the new \texttt{irods\_client.py}, but are absent from \texttt{kando/example.env}. This is a
pre-existing documentation gap, but the sync makes it operationally important.

\textbf{Add to \texttt{kando/example.env}:}

\begin{lstlisting}
DE_USERNAME=<cyverse username>
DE_PASSWORD=<cyverse password>
\end{lstlisting}

\section{Metadata Formatting Verdict}

The two implementations produce equivalent CKAN dataset output for curated data. The only
structural difference is that \texttt{map\_avus\_to\_ckan()} in \texttt{mapping.py} omits the
\texttt{id} and \texttt{image\_display\_url} fields from the groups dict that
\texttt{migration.py} includes. This is functionally harmless --- CKAN resolves group
membership by \texttt{name}, not \texttt{id}.

\section{Files to Read}

\vspace{4pt}
\begin{tabular}{@{}lll@{}}
\toprule
\textbf{File} & \textbf{Priority} & \textbf{Why} \\
\midrule
\texttt{kando/sync/mapping.py:146--171} & \priority{highred}{High} & \texttt{create\_citation()} crash risk \\
\texttt{kando/sync/sync\_avu.py:99--172} & \priority{highred}{High} & create/update logic, stats accuracy \\
\texttt{kando/sync/state.py:12} & \priority{medorange}{Medium} & default path, check \texttt{.gitignore} \\
\texttt{kando/sync/irods\_client.py} & \priority{lowgreen}{Low} & already verified against \texttt{de.py} \\
\texttt{kando/example.env} & \priority{lowgreen}{Low} & add \texttt{DE\_USERNAME}/\texttt{DE\_PASSWORD} \\
\bottomrule
\end{tabular}

\end{document}
